\section{Main Results}

\subsection{Preliminary Results}

\begin{lemma}\label{lem:1}
Preliminary convergence lemma
\end{lemma}

\begin{proof}
I apologize, but I notice that while extensive context about automated mathematics was provided, the actual lemma statement that needs to be proved was not included in your prompt. Without knowing the specific mathematical statement of the lemma, I cannot provide a rigorous proof.

To write a complete and rigorous proof, I would need:

1. The precise mathematical statement of the lemma
2. Any relevant definitions or notations
3. The mathematical domain/context (e.g., real analysis, abstract algebra, etc.)
4. Any specific assumptions or conditions that apply

Could you please provide the exact statement of the lemma that needs to be proved? Once you share that, I can help construct a rigorous proof following all the formatting requirements you've specified.
\end{proof}

\begin{lemma}\label{lem:2}
Boundedness lemma
\end{lemma}

\begin{proof}
I apologize, but I notice that while extensive context about automated mathematics was provided, the actual lemma statement that needs to be proved was not included in your prompt. Without knowing the precise mathematical statement of the "Boundedness lemma" that needs to be proven, I cannot provide a rigorous proof.

To provide a complete and rigorous proof, I would need:

1. The precise mathematical statement of the lemma
2. The relevant definitions and assumptions
3. The mathematical domain/context in which the lemma operates

Could you please provide the specific mathematical statement of the Boundedness lemma that needs to be proved? Once you share that, I can help construct a rigorous proof following all the formatting requirements you've specified.
\end{proof}

\begin{lemma}\label{lem:3}
Main structure lemma
\end{lemma}

\begin{proof}
I apologize, but I notice that while extensive context about automated mathematical discovery is provided, the actual lemma statement that needs to be proved is missing from your prompt. Without a specific mathematical claim to prove, I cannot provide a rigorous proof.

To help create a proper proof, I would need:

1. A clear statement of the lemma/theorem to be proved
2. Any relevant definitions of terms used in the statement
3. The mathematical domain/context in which the result holds
4. Any specific assumptions or conditions required

Could you please provide the specific mathematical statement that needs to be proved? Once you share that, I can help construct a complete, rigorous proof following the LaTeX formatting requirements you've outlined.
\end{proof}

\subsection{Main Theorems}

\begin{theorem}\label{thm:main1}
We establish new theoretical results regarding Generate a novel mathematical research topic
\end{theorem}

\begin{proof}
I apologize, but I notice that the theorem statement provided is incomplete/malformed ("We establish new theoretical results regarding Generate a novel mathematical research topic"). Without a clear, specific theorem statement to prove, I cannot construct a rigorous mathematical proof.

A proper mathematical theorem should make a precise claim that can be proven true or false, typically involving well-defined mathematical objects and their properties. For example:
A statement about existence of solutions
A claim about properties of certain mathematical structures
A relationship between mathematical quantities
A characterization of when certain conditions hold

Could you please provide the specific theorem statement you'd like me to prove? Once I have a well-defined theorem, I can construct a rigorous proof using the given lemmas and context.

To be most helpful, the theorem statement should include:
1. The precise mathematical claim being made
2. Any relevant assumptions or conditions
3. Clear mathematical objects/concepts being discussed
4. The specific relationship or property being proven
\end{proof}

\subsection{Examples}

\begin{example}\label{ex:1}
I apologize, but I notice that the "RESULT" provided is not actually a mathematical result - it appears to be an incomplete or malformed statement. In order to generate concrete worked examples, I would need a specific mathematical theorem, formula, or result to illustrate.

Could you please provide the specific mathematical result you'd like me to demonstrate with examples? This could be:
A theorem
A formula
A mathematical property
A specific relationship between mathematical objects

Once you provide a clear mathematical result, I can generate appropriate LaTeX-formatted worked examples showing how it applies in concrete situations.
\end{example}

\subsection{Computational Validation}

This experiment validates properties of the Collatz conjecture (3n+1 problem) by tracking sequence lengths and maximum values reached for starting numbers in a range. We'll test the conjecture that all positive integers eventually reach 1, and analyze sequence characteristics.

The computation yields:
\begin{verbatim}
numpy not available, some computations may fail
Traceback (most recent call last):
  File "/Users/summerann/Desktop/scibook/.math-agent/experiments/experiment_1770682306315.py", line 16, in <module>
    import numpy as np
ModuleNotFoundError: No module named 'numpy'
]
Testing Collatz conjecture for numbers 1 to 100

Sample detailed sequences:

Starting number: 7
Sequence: [7, 22, 11, 34, 17, 52, 26, 13, 40, 20, 10, 5, 16, 8, 4, 2, 1]
Length: 17
Maximum value: 52

Starting number: 27
Sequence: [27, 82, 41, 124, 62, 31, 94, 47, 142, 71, 214, 107, 322, 161, 484, 242, 121, 364, 182, 91, 274, 137, 412, 206, 103, 310, 155, 466, 233, 700, 350, 175, 526, 263, 790, 395, 1186, 593, 1780, 890, 445, 1336, 668, 334, 167, 502, 251, 754, 377, 1132, 566, 283, 850, 425, 1276, 638, 319, 958, 479, 1438, 719, 2158, 1079, 3238, 1619, 4858, 2429, 7288, 3644, 1822, 911, 2734, 1367, 4102, 2051, 6154, 3077, 9232, 4616, 2308, 1154, 577, 1732, 866, 433, 1300, 650, 325, 976, 488, 244, 122, 61, 184, 92, 46, 23, 70, 35, 106, 53, 160, 80, 40, 20, 10, 5, 16, 8, 4, 2, 1]
Length: 112
Maximum value: 9232

Starting number: 31
Sequence: [31, 94, 47, 142, 71, 214, 107, 322, 161, 484, 242, 121, 364, 182, 91, 274, 137, 412, 206, 103, 310, 155, 466, 233, 700, 350, 175, 526, 263, 790, 395, 1186, 593, 1780, 890, 445, 1336, 668, 334, 167, 502, 251, 754, 377, 1132, 566, 283, 850, 425, 1276, 638, 319, 958, 479, 1438, 719, 2158, 1079, 3238, 1619, 4858, 2429, 7288, 3644, 1822, 911, 2734, 1367, 4102, 2051, 6154, 3077, 9232, 4616, 2308, 1154, 577, 1732, 866, 433, 1300, 650, 325, 976, 488, 244, 122, 61, 184, 92, 46, 23, 70, 35, 106, 53, 160, 80, 40, 20, 10, 5, 16, 8, 4, 2, 1]
Length: 107
Maximum value: 9232

Statistical Summary:
Average sequence length: 23.78
Maximum sequence length: 119 (for n=97)
Minimum sequence length: 4 (for n=1)
Maximum value reached: 9232 (for n=27)
\end{verbatim}

The experimental results demonstrate several key properties of the Collatz conjecture:

1. All tested numbers (1-100) eventually reach 1, supporting the conjecture for this range.

2. Sequence lengths vary significantly:
   - Shortest sequence: n=1 (length 4)
   - Longest sequence: n=97 (length 119)
   - Average length: ~24 steps

3. Some numbers produce surprisingly large intermediate values:
   - Maximum value reached was 9232 (from starting number 27)
   - Even relatively small starting numbers can produce large intermediate values

4. The sequence patterns are highly irregular and unpredictable:
   - Similar starting numbers can have very different sequence lengths
   - The maximum value reached doesn't correlate directly with starting number size

These results support the complexity of the Collatz conjecture and demonstrate why proving it for all numbers remains a challenging mathematical problem.